By carbon leakage, we mean the reallocation of economic activity — and, as a consequence, of emissions — from one economic unit to another in response to environmental regulation. When introducing such policies, one of policymakers' primary concerns is the international margin of leakage — the reallocation of emission-intensive production abroad — both because it hurts national output and employment and because it undermines the effectiveness of the policy. At its core, however, the mechanism is more general: even within a jurisdiction, environmental regulation can create winners and losers. Documenting under what circumstances and to what degree leakage occurs, along either margin, is an essential ingredient for understanding the effect of environmental regulation, and, in particular, carbon pricing, on aggregate emissions.

In this paper, we jointly test for the domestic and international leakage margins in the context of the European Union Emissions Trading System (EU ETS) in Belgium. We find a null result on both.

The Belgian context is well-suited to this test. As an EU member state, Belgium has been subject to a common European carbon price since 2005, with each regulated installation's annual emissions and allowance allocation recorded administratively over the entire period. Unlike any other EU member state, Belgium also offers firm-level visibility into how firms source their inputs: a near-universal panel of domestic firm-to-firm transactions on one side, and a firm-by-product-by-source-country panel of imports on the other, both spanning the period before and after the carbon price was introduced. This allows us to focus on the buyers for whom reorganizing the sourcing decision would be most profitable — those whose input costs are most directly hit by the regulated suppliers under the heaviest carbon pressure — and ask whether they do so on either the intensive or extensive margin, in the short and medium run.

Our analysis provides three sets of results. First, we document that putting a price on carbon emissions increases prices in the manufacturing sector in Belgium, and more so in sectors heavily exposed to the policy.

Second, we show that [domestic margin].

Third, we document [international margin].

We then proceed to offer possible explanations that can account for the empirical findings above. [RE-WRITE WHAT FOLLOWS, BUT THE SUBSTANCE IS THERE]. The most natural --- that the cost shock at the typical buyer's input bill is too small to overcome any reasonable cost of switching suppliers --- is consistent with the data on average but is harder to defend among the most exposed buyers, where we still find no reallocation. We rule out one mechanism directly: that long-standing incumbency in the regulated supplier relationship is the binding friction. If long-tenured incumbents were sheltered by the relationship itself, the post-2015 share loss for the most-exposed regulated supplier would be smallest where that supplier is also the buyer's longest-tenured incumbent. The pattern is the opposite. Three other candidates remain consistent with the data we have, and discriminating among them requires data we do not: that competing suppliers move prices together and so compress the relative-price signal that would induce buyer substitution, that pass-through from the carbon shock onto regulated suppliers' prices is only partial, and that the ``unregulated'' alternatives are themselves partially regulated through their own upstream regulated inputs.

Our findings contribute to two strands of the literature. The first is the cross-border leakage literature, which has measured the response of import flows to unilateral environmental regulation using aggregate trade data \citep{aichelefelbermayr2015,naegelezaklan2019,satodechezlepretre2015} and, most recently, at the firm level using French customs data \citep{CosterMejeanDiGiovanni2024}. We run the same firm-level design on Belgian customs data over the same period, and find that the substantial reallocation toward non-EU sources documented in France does not appear in Belgium. The two economies face the same EU-wide carbon-pricing regime; why their import responses diverge is an open question we do not tackle. The second is the firm-level literature on within-country EU ETS effects \citep{martin2014,dechezlepretreetal2023,colmeretal2023,kanzig2023unequal}, which documents that regulated firms reduced emissions without contracting output relative to unregulated controls. This literature has used the average effect on regulated firms; it has not exploited firm-level heterogeneity in carbon-price exposure within the regulated group. A firm-level null on the average treated effect is silent on whether more-exposed regulated firms lose share to less-exposed regulated firms in the same product market. That margin --- within the regulated group, between firms heterogeneous in exposure --- is where reallocation should be hardest to miss if it is operating. We test it directly and find a null there as well. We also focus on the buyer's response rather than the regulated firm's, a margin where reorganization could appear even when the regulated firms themselves do not visibly contract.
